\chapter{1937:Norman Haworth and Paul Karrer}

\label{chap:chemistry1937}

The Nobel Prize in Chemistry for the year 1937 was awarded jointly to Norman Haworth and Paul Karrer.

\textbf{Motivation:}

for his investigations on carbohydrates and vitamin C (one half, Norman Haworth) and for his investigations on carotenoids, flavins and vitamins A and B2 (the other half, Paul Karrer)

\section{Norman Haworth}

\subsection*{Early Life and Education}

Walter Norman Haworth was born on 1883-03-19 in Coventry, England.
He was a British chemist who shared the 1937 Nobel Prize in Chemistry with Paul Karrer for his investigations on carbohydrates and vitamin C.

\subsection*{Career and Major Research}

Haworth studied at the University of Manchester and received his doctorate at the University of Göttingen in 1910 under Otto Wallach, working on terpenes.
He then worked with Thomas Purdie at the University of St Andrews, applying the methylation method to the sugars.
He was professor at Armstrong College in Newcastle upon Tyne from 1920 to 1925, and from 1925 he held the Mason Chair of Chemistry at the University of Birmingham until his retirement in 1948.
He established that the sugars exist in ring forms and introduced the formulas that bear his name, determining the ring size and configuration of many sugars and their glycosides.
In 1933 his group, working with Edmund Hirst, achieved the total synthesis of vitamin C (ascorbic acid), confirming the structure of the substance that had been isolated by Albert Szent-Györgyi.

\subsection*{Nobel Prize-winning Work}

The work for which Norman Haworth was awarded the Nobel Prize in Chemistry in 1937 was described by the Nobel Committee as: "for his investigations on carbohydrates and vitamin C".

Haworth's investigations established the ring structures of the sugars and their derivatives, and they culminated in the first total synthesis of vitamin C.
The synthetic route he developed also made possible the industrial production of the vitamin.

\subsection*{Significance and Impact}

Haworth's work transformed the chemistry of carbohydrates into a field in which the configuration and ring structure of every sugar could be specified.
The synthesis of vitamin C was an early triumph of the chemistry of natural products and made the vitamin available for medicine and nutrition.

\subsection*{Later Career and Honors}

Haworth was elected a Fellow of the Royal Society in 1928 and was knighted in 1947.
He received the Davy Medal of the Royal Society in 1934.
He died on 1950-03-19 in Birmingham, England.

\subsection*{Legacy}

The legacy of Norman Haworth endures through the lasting impact of his scientific contributions.
The Haworth formulas and his methods of sugar chemistry are still taught and used throughout carbohydrate chemistry.

\subsection*{Modern Developments}

Haworth's elucidation of carbohydrate ring structures founded the modern chemistry of sugars and polysaccharides, underpinning glycobiology, the food industry, and the production of cellulose-based materials. His synthesis of vitamin C enabled its large-scale industrial production, and vitamin C supplements and fortified foods remain central to nutrition worldwide.

\subsection*{Selected Publications}

\begin{itemize}

\item Haworth, W. N. (1929). "The Constitution of Sugars." London: Edward Arnold.

\item Haworth, W. N., \& Hirst, E. L. (1933). "Synthesis of Ascorbic Acid." Journal of the Chemical Society.

\end{itemize}

\clearpage

\section{Paul Karrer}

\subsection*{Early Life and Education}

Paul Karrer was born on 1889-04-21 in Moscow, Russia, to Swiss parents. His family returned to Switzerland in 1892, and he was educated in Aarau and at the University of Zurich, where he studied under Alfred Werner and received his doctorate in 1911.
He was a Swiss organic chemist who shared the 1937 Nobel Prize in Chemistry with Norman Haworth for his investigations on carbohydrates and vitamin C.

\subsection*{Career and Major Research}

Karrer studied at the University of Zürich, where he received his doctorate in 1911 under Alfred Werner.
He then worked for several years with Paul Ehrlich in Frankfurt am Main, where he studied organic arsenic compounds and their chemotherapeutic action.
In 1918 he returned to Zürich, and in 1919 he became professor of organic chemistry at the university, succeeding his teacher Alfred Werner; he remained there for the rest of his career.
His research concerned the plant pigments known as carotenoids; he and his co-workers determined the constitution of the carotenes and showed that the yellow pigment beta-carotene is the precursor from which vitamin A is formed in the body.
He also studied the flavins and contributed to the chemistry of the coenzymes involved in cellular oxidation.

\subsection*{Nobel Prize-winning Work}

The work for which Paul Karrer was awarded the Nobel Prize in Chemistry in 1937 was described by the Nobel Committee as: "for his investigations on carotenoids, flavins and vitamins A and B2".

Karrer's elucidation of the constitution of the carotenoid pigments, and his proof that carotene is converted into vitamin A in the body, connected the plant pigments to the chemistry of the vitamins.

\subsection*{Significance and Impact}

Karrer's work established the chemistry of the carotenoids and clarified the relationship between carotene and vitamin A, which is the basis of the vitamin A activity of vegetables.
His textbook of organic chemistry, first published in 1928, was used by students throughout the world.

\subsection*{Later Career and Honors}

Karrer remained at the University of Zürich until his retirement in 1959.
He held honorary doctorates from several universities and was a member of several academies.
He died on 1971-06-18 in Zürich, Switzerland.

\subsection*{Legacy}

The legacy of Paul Karrer endures through the lasting impact of his scientific contributions.
The chemistry of the carotenoids and of vitamin A that he established continues to be built upon by researchers across the chemical and biological sciences.

\subsection*{Modern Developments}

Karrer's investigations of carotenoids, flavins, and vitamins A and B2 established the structural and synthetic chemistry of vitamins. His work enabled the industrial production of vitamin A, used today to fortify foods and prevent blindness, and established flavins as components of cellular redox coenzymes. Vitamin and carotenoid research continues to inform nutrition and medicine.

\subsection*{Selected Publications}

\begin{itemize}

\item Karrer, P. (1928). "Lehrbuch der organischen Chemie." Leipzig: Georg Thieme.

\item Karrer, P. (1937). Nobel Lecture on the carotenoids and vitamins, Stockholm.

\end{itemize}

\clearpage

\section*{Chapter Summary}

The 1937 prize was shared by Norman Haworth for his investigations on carbohydrates and vitamin C, and by Paul Karrer for his investigations on carotenoids and vitamins. Together their work established the structure and synthesis of the vitamins and carbohydrates.
